\documentclass[../../main.tex]{subfiles}
\begin{document}

{\bf [解答]}

\paragraph{(1)}

$4$回操作した時，$x=\pm1,\pm3$にある確率は明らかに$0$であることに注意する．
又，$x=\pm2$にいる確率は対称性から等しく，例えば表が$3$回出る時で
\begin{align}
  \dfrac{{}_4C_1}{2^4}=\dfrac{1}{4}\label{eq:1}
\end{align}
$x=\pm4$にいる確率は対称性から等しく，例えば表が$4$回出る時で
\begin{align}
  \dfrac{{}_4C_0}{2^4}=\dfrac{1}{16}\label{eq:2}
\end{align}
\cref{eq:1}，\cref{eq:2}から，余事象より$x=0$にいる確率は
\begin{align}
  1-2\left(\dfrac{1}{4}+\dfrac{1}{16}\right)=\dfrac{3}{8}
\end{align}
である．以上をまとめて
\begin{align}
  \begin{cases}
    x=\pm1,\pm3 & 0     \\
    x=0                & 3/8   \\
    x=\pm2          & 1/4   \\
    x=\pm4          & 1/16
  \end{cases}
\end{align}
である．$\cdots$(答)

\paragraph{(2)}

$x=n-2$にいるのは表が$n-1$回出た時で$\dfrac{{}_nC_{n-1}}{2^n}=\dfrac{n}{2^n}$．
$\cdots$(答)
また$x=n-4$にいるのは表が$n-2$回出た時で$\dfrac{{}_nC_{n-2}}{2^n}=\dfrac{n(n-1)}{2^{n+1}}$．
$\cdots$(答)
\newpage
\end{document}
